\documentclass{beamer}

\mode<presentation>
\usepackage{amsmath,amssymb,mathtools}
\usepackage{textcomp}
\usepackage{gensymb}
\usepackage{adjustbox}
\usepackage{subcaption}
\usepackage{enumitem}
\usepackage[utf8]{inputenc}
\usepackage{amssymb}
\usepackage{enumitem}
\setlist{nosep} % optional: removes vertical gaps
\setlist[enumerate]{label=\arabic*)} % custom numbering if you want

\usepackage{multicol}
\usepackage{listings}
\usepackage{url}
\usepackage{graphicx} % <-- needed for images
\def\UrlBreaks{\do\/\do-}

\usetheme{Boadilla}
\usecolortheme{lily}
\setbeamertemplate{footline}{
  \leavevmode%
  \hbox{%
  \begin{beamercolorbox}[wd=\paperwidth,ht=2ex,dp=1ex,right]{author in head/foot}%
    \insertframenumber{} / \inserttotalframenumber\hspace*{2ex}
  \end{beamercolorbox}}%
  \vskip0pt%
}
\setbeamertemplate{navigation symbols}{}

\lstset{
  frame=single,
  breaklines=true,
  columns=fullflexible,
  basicstyle=\ttfamily\small % Adjusted for better readability
}

\numberwithin{equation}{section}

% ---- your macros ----
\providecommand{\nCr}[2]{\,^{#1}{#2}}
\providecommand{\nPr}[2]{\,^{#1}P_{#2}}
\providecommand{\mbf}{\mathbf}
\providecommand{\pr}[1]{\ensuremath{\Pr\left(#1\right)}}
\providecommand{\qfunc}[1]{\ensuremath{Q\left(#1\right)}}
\providecommand{\sbrak}[1]{\ensuremath{{}\left[#1\right]}}
\providecommand{\lsbrak}[1]{\ensuremath{{}\left[#1\right.}}
\providecommand{\rsbrak}[1]{\ensuremath{\left.#1\right]}}
\providecommand{\brak}[1]{\ensuremath{\left(#1\right)}}
\providecommand{\lbrak}[1]{\ensuremath{\left(#1\right.}}
\providecommand{\rbrak}[1]{\ensuremath{\left.#1\right)}}
\providecommand{\cbrak}[1]{\ensuremath{\left\{#1\right\}}}
\providecommand{\lcbrak}[1]{\ensuremath{\left\{#1\right.}}
\providecommand{\rcbrak}[1]{\ensuremath{\left.#1\right\}}}
\theoremstyle{remark}


\providecommand{\abs}[1]{\left\vert#1\right\vert}
\providecommand{\res}[1]{\Res\displaylimits_{#1}}
\providecommand{\norm}[1]{\lVert#1\rVert}
\providecommand{\mtx}[1]{\mathbf{#1}}
\providecommand{\mean}[1]{E\left[ #1 \right]}
\providecommand{\fourier}{\overset{\mathcal{F}}{ \rightleftharpoons}}

\providecommand{\dec}[2]{\ensuremath{\overset{#1}{\underset{#2}{\gtrless}}}}

\usepackage{gvv}
% ---------------------

\title{Matgeo Presentation - Problem 8.2.21}
\author{EE25BTECH11054 - S.Harsha Vardhan}

\begin{document}

%---------------- Title Page ----------------
\begin{frame}
  \titlepage
\end{frame}

%---------------- Problem Statement ----------------
\begin{frame}{Problem Statement}
Find the equation of the conic that satisfies the given conditions:
\begin{itemize}
    \item vertex (0,0)
    \item passing through (5,2)
    \item symmetric with respect to Y axis.
\end{itemize}
\end{frame}

%---------------- Solution ----------------
\begin{frame}{Solution: General Equation}
The general conic equation is:
\begin{align}
 g(\vec{x}) = \vec{x}^\top \vec{V}\vec{x} + 2\vec{u}^\top \vec{x} + f = 0 \tag{1}
\end{align}
For a parabola with vertex at the origin and axis along the Y-axis, the focus is $\vec{F} = \myvec{0 \\ p}$ and the directrix is $y = -p$. This gives $\vec{n} = \myvec{0 \\ 1}$ and $c = -p$.
\end{frame}

\begin{frame}{Solution: Deriving Parameters}
\begin{align}
	\vec{V} &= \|\vec{n}\|^2 \vec{I} - \vec{n}\vec{n}^\top 
	= \myvec{1 & 0 \\ 0 & 1} - \myvec{0 \\ 1}\myvec{0 & 1} 
	= \myvec{1 & 0 \\ 0 & 0} \tag{2}
\end{align}
\begin{align}
	\vec{u} &= c\vec{n} - \vec{F} 
	= (-p)\myvec{0 \\ 1} - \myvec{0 \\ p} 
	= \myvec{0 \\ -2p} \tag{3}
\end{align}
\begin{align}
	f &= \|\vec{F}\|^2 - c^2 = p^2 - (-p)^2 = 0 \tag{4}
\end{align}
\end{frame}

\begin{frame}{Solution: Equation of Parabola}
Substituting these into the general conic equation (1):
\begin{align}
	\vec{x}^\top \myvec{1 & 0 \\ 0 & 0}\vec{x} 
	+ 2\myvec{0 & -2p}^\top\vec{x} + 0 &= 0 \tag{5} \\
	x^2 - 4py &= 0 \tag{6}
\end{align}
The curve passes through the point (5,2).
\begin{align}
	(5)^2 - 4p(2) &= 0 \tag{7} \\
	25 - 8p &= 0 \tag{8} \\
	p &= \frac{25}{8} \tag{9}
\end{align}
\end{frame}

\begin{frame}{Solution: Final Equation}
The final parameters are:
\begin{align}
	\vec{V} = \myvec{1 & 0 \\ 0 & 0}, \quad
	\vec{u} = \myvec{0 \\ -\frac{25}{4}}, \quad
	f = 0 \tag{10}
\end{align}
The equation of the parabola is:
\begin{align}
	x^2 &= 4\left(\frac{25}{8}\right)y \tag{11} \\
	x^2 &= \frac{25}{2}y \tag{12}
\end{align}
\end{frame}

\begin{frame}{Plot}
    \begin{figure}[h]
    \centering
    \includegraphics[width=0.8\columnwidth]{../figs/Figure_10.png}
    \caption*{}
    \label{fig:placeholder}
    \end{figure}
\end{frame}

\section{C Code}
\begin{frame}[fragile]{C Code: parabola\_calculator.c}
\begin{lstlisting}[language=C]
double calculate_y(double x) {
    double y = (2.0 / 25.0) * x * x;
    return y;
}
\end{lstlisting}
\end{frame}

\section{Python Code}
\begin{frame}[fragile]{Python Code: main\_caller.py}
\begin{lstlisting}[language=Python]
import ctypes
import os

lib_path = os.path.join(os.getcwd(), 'libparabola.so')
parabola_lib = ctypes.CDLL(lib_path)

parabola_lib.calculate_y.argtypes = [ctypes.c_double]
parabola_lib.calculate_y.restype = ctypes.c_double

x_test = 5.0
y_result = parabola_lib.calculate_y(x_test)

print(f"Calling C function with x = {x_test}")
print(f"The calculated y value is: {y_result}")
\end{lstlisting}
\end{frame}

\begin{frame}[fragile]{Python Code: plotter.py}
\begin{lstlisting}[language=Python]
import numpy as np
import matplotlib.pyplot as plt

x_vals = np.linspace(-8, 8, 400)
y_vals = (2.0 / 25.0) * x_vals**2

plt.figure(figsize=(8, 6))

plt.plot(x_vals, y_vals, label=r'$y = \frac{2}{25}x^2$', color='blue')

plt.scatter([0], [0], color='red', zorder=5)
plt.annotate('Vertex (0,0)', xy=(0, 0), xytext=(0.5, 1))

plt.scatter([5], [2], color='green', zorder=5)
plt.annotate('Point (5,2)', xy=(5, 2), xytext=(2.5, 3))

plt.title('Parabola with Vertex at (0,0)')
plt.xlabel('x-axis')
plt.ylabel('y-axis')
plt.grid(True)
plt.axis('equal')
plt.legend()

plt.show()
\end{lstlisting}
\end{frame}

\end{document}